%% Research Note of the AAS (RNAAS): G 249-11 (TIC 417732194).
%% Compile on Overleaf with the AAS template "AASTeX Template for submissions to AAS
%% Journals (ApJ-AJ-ApJS-ApJL-PSJ-RNAAS)", which provides aastex701.cls: upload this
%% file, references.bib and figure1.pdf to that project. RNAAS limits: 1,500 words in
%% all (title, headers, caption and references included), 150 for the abstract, and
%% one figure; no line numbers. Every number comes from results/g249-11 in
%% https://github.com/Comdex4/tess-transit-hunter (see README.md in this folder).
\documentclass{aastex701}

%% Collaboration authors have no initials; fixed citation text avoids a stray space.
\defcitealias{Gaia2023}{Gaia Collaboration et al.\ 2023}
\defcitealias{Lightkurve2018}{Lightkurve Collaboration et al.\ 2018}
\defcitealias{Astropy2022}{Astropy Collaboration et al.\ 2022}

\begin{document}

\title{A Candidate Transiting Super-Earth around the Nearby M4.5 Dwarf G 249-11}

%% Add the e-mail address the journal should show.
\author{Connor Rice}
\affiliation{Independent Researcher, USA}
\email[show]{EMAIL}

\begin{abstract}
I report a transiting planet candidate around G~249-11 (TIC~417732194), an M4.5 dwarf
29~pc away. Transit Least Squares finds a 5.3075-day signal in TESS 2-minute photometry
of sectors 19, 59 and 60 (signal detection efficiency 21.2), and the same period in
sectors 19, 73 and 86, without the discovery data (SDE 9.8). The signal passes all 17
flux-level tests of the LEO-Vetter automated vetter. TRICERATOPS gives a false-positive
probability of 0.10, mostly from an unresolved bound companion. A transit fit gives a
depth of 3100~ppm, a duration of 0.86~h and a radius of $1.6^{+0.4}_{-0.15}\,R_\oplus$;
the planet would receive 4.8 times Earth's insolation. The star is bright in the
infrared ($K_s=9.7$), and the radial-velocity semi-amplitude would be
2--6~m~s$^{-1}$, so the candidate is accessible to follow-up.
\end{abstract}

\keywords{\uat{Exoplanet detection methods}{489} --- \uat{Transit photometry}{1709} ---
\uat{M dwarf stars}{982}}

\section{Introduction}

Small planets transiting nearby M dwarfs are the most accessible rocky planets for mass
and atmosphere measurements. I searched the TESS \citep{Ricker2015} 2-minute light curves
of up to 1000 M dwarfs per year with a box least-squares search \citep{Kovacs2002} and
found a 5.3-day signal on G~249-11 in sectors 59 and 60. The star has no TESS Object of
Interest or community TOI and is not among the 172 M-dwarf candidates of
\citet{Kunimoto2025}. Here I detect and vet the signal with published tools.

\section{Data}

G~249-11 (dM4.5; \citealt{Hejazi2022}) has a Gaia DR3 parallax of $34.378\pm0.017$~mas
(29.1~pc) and RUWE of 1.20 \citepalias{Gaia2023}. TIC~8 \citep{Stassun2019} gives
$R_\star=0.266\pm0.008\,R_\odot$, $M_\star=0.238\pm0.020\,M_\odot$ and
$T_{\rm eff}=3180\pm157$~K. TESS observed it at 2-minute cadence in sectors 19, 59 and 60,
and in 200-s full-frame images in sectors 73 and 86. I used the SPOC PDCSAP light curves
\citep{Jenkins2016} and the QLP detrended light curves \citep[cadences with QUALITY
$=0$;][]{Huang2020}, downloaded with Lightkurve \citepalias{Lightkurve2018}, and detrended
each sector with wotan's biweight filter \citep[0.75-day window;][]{Hippke2019wotan}.

\section{Detection}

Transit Least Squares \citep[TLS;][]{Hippke2019tls}, searching 0.5--20 days with the
star's limb darkening \citep{Claret2017}, finds its highest peak at 5.3097~d (signal
detection efficiency, SDE, 11.8) in sectors 59 and 60 alone, and at
$5.30748\pm0.00017$~d (SDE 21.2) in all three 2-minute sectors, where the next peak, its
$P/3$ alias, reaches 12.9. Without the discovery data, in sectors 19, 73 and 86, the
second-highest peak is at 5.30742~d (SDE 9.8). The highest (9.92~d, SDE 10.6) rests on
one dip in sector 19 and three in sector 86: at its predicted transits in sector 59 the
flux drops by only $340\pm450$~ppm, and its depth differs between sectors ($\chi^2=15.6$
for 4 degrees of freedom, $p=0.004$). At the candidate's period every sector shows a dip,
of consistent depth ($\chi^2=5.8$, $p=0.21$). All five sectors together, with 18
transits, give $5.30742\pm0.00014$~d at SDE 26.0. Above SDE 9.1, the TLS false-alarm
probability is below $10^{-4}$ \citep{Hippke2019tls}.

\begin{figure}[t!]
\plotone{figure1.pdf}
\caption{TESS photometry of G~249-11 folded at the fitted ephemeris: (a)--(c) one panel
per epoch, with the individual measurements (gray), 15-minute bins (blue) and the
maximum-posterior transit model (orange); (d) the 13 individual 2-minute transits in
30-minute bins, each with the same model.
\label{fig:transits}}
\end{figure}

\section{Vetting}

LEO-Vetter \citep{Kunimoto2025} tests a signal against 13 kinds of false alarm (noise,
systematics, variability) and 4 kinds of astrophysical false positive (eclipsing
binaries). With its default thresholds, the signal in the 2-minute data passes all 17
tests: multiple event statistic 9.2, consistent depths and S/N from transit to transit,
a transit model strongly preferred over a flat line ($\Delta{\rm AIC}=-114$), odd and even
transits within $0.3\sigma$, and no significant secondary eclipse. With all five sectors
(18 transits) it again passes all 17, with a multiple event statistic of 10.6.

TRICERATOPS \citep{Giacalone2021}, given the folded 2-minute light curve, the SPOC
apertures, the 144 TIC stars within 10 pixels and the Gaia DR3 field population, gives a
false-positive probability FPP $=0.107\pm0.005$ and a nearby false-positive probability
NFPP $=0.0039\pm0.0001$ (mean and scatter of 10 runs). The FPP comes almost entirely from
a planet transiting an unresolved bound companion (0.10), and the NFPP from the three
neighbors bright enough to cause the dip (20--43\arcsec\ away, needing 15--65\% eclipses).
Both exceed the thresholds for statistical validation (FPP $<0.015$, NFPP $<0.001$), so
the candidate is not validated, although the NFPP is far below the 0.1 that marks a
likely nearby false positive. High-resolution imaging and seeing-limited photometry of a
transit would settle both.

\section{Properties}

A fit of a batman model \citep{Kreidberg2015} to the 2-minute data with emcee
\citep{ForemanMackey2013} gives $P=5.307419\pm0.000010$~d,
$T_0={\rm BJD_{TDB}}\,2459922.9291\pm0.0012$, a depth of $3090^{+1650}_{-500}$~ppm, a
duration of $0.86^{+0.11}_{-0.08}$~h and $R_p=1.6^{+0.4}_{-0.15}\,R_\oplus$. The
semi-major axis would be 0.037~AU, with 4.8 times Earth's insolation and an equilibrium
temperature of 410~K (zero albedo). For 2.5--6~$M_\oplus$ the radial-velocity
semi-amplitude would be 2.4--5.7~m~s$^{-1}$. Figure~\ref{fig:transits} shows the transit
in every epoch.

\section{Conclusion}

G~249-11 hosts a candidate super-Earth on a 5.3-day orbit that passes automated vetting
but is not yet statistically validated. Ground-based transit photometry,
high-resolution imaging and radial velocities are needed to confirm it. The analysis
scripts and their outputs are at \url{https://github.com/Comdex4/tess-transit-hunter}.

\begin{acknowledgments}
This Note uses TESS data, funded by NASA's Science Mission Directorate and obtained from
MAST at STScI
(\dataset[doi:10.17909/t9-nmc8-f686]{https://doi.org/10.17909/t9-nmc8-f686},
\dataset[doi:10.17909/t9-r086-e880]{https://doi.org/10.17909/t9-r086-e880}).
I used Claude, an AI assistant made by Anthropic, to write and run the analysis code at
my direction and to draft and edit this Note. I checked the code, the results and the
text, and I take full responsibility for the content.
\end{acknowledgments}

\facilities{TESS, Gaia}

\software{Lightkurve \citepalias{Lightkurve2018}, wotan \citep{Hippke2019wotan},
TLS \citep{Hippke2019tls}, LEO-Vetter \citep{Kunimoto2025},
TRICERATOPS \citep{Giacalone2021}, batman \citep{Kreidberg2015},
emcee \citep{ForemanMackey2013}, Astropy \citepalias{Astropy2022}}

\bibliography{references}{}
\bibliographystyle{aasjournalv7}

\end{document}
